\begin{statement}{C.2.UN.51}{Proved}{The mixed-cycle criterion and compact occurrence with Hopf components}
Specify an oriented cycle of surfaces, each of which is either
a diagonal primary or secondary Hopf surface of algebraic dimension
zero, marked by its two axis curves, or
\[
                  R_i=\PP(\OO_{E_i}\oplus L_i),
                  \qquad E_i\ \hbox{elliptic},\quad
                         L_i\in\operatorname{Pic}^0(E_i),
                                                        \tag{C.MX.1}
\]
marked by its zero and infinity sections. Specify the elliptic
isomorphism joining the outgoing curve to the next incoming curve,
and a trivialization of their tensor product of normal lines.
The associated proper normal-crossing plumbing has two-torus
general fibre exactly when the cycle product $M$ of its actual
rank-three integral complement lattices is the identity.
Every actual reduced smooth-total two-torus smoothing of the
specified support satisfies this criterion.

For every such cycle containing a Hopf component, $M=I$ is also
equivalent to occurrence of the entire specified reduced support
in a smooth compact threefold over $\PP^1$, with two-torus
general fibre and Euler number zero for every reduced fibre
support. This retains all component surfaces and every elliptic
gluing map. The construction chooses the normal-product constants.
A compatible cycle containing Hopf components contains at least
two of them.

The local criterion also includes all-ruled cycles and the
identification of the two marked sections of one ruled surface.
In that case the exact punctured period lattice has the form
\[
                 \Lambda+\Z\bigl(b_0+v_0\log t\bigr),
           \qquad \Lambda\simeq\Z^3\subset\C^2,          \tag{C.MX.2}
\]
where the projection of $b_0$ to the elliptic quotient of the
common complement is the actual total translation of the base
gluing maps. It is retained in this formula.
\end{statement}

\begin{proof}
\par\medskip\noindent\textbf{The complements and their actual line-bundle maps.}
For a ruled component, removing the two sections gives the
nonzero line bundle $L_i^\times$. Choose an origin on $E_i$
and a rigidification of $L_i$ there. Since $\deg L_i=0$, it
has constant factors of automorphy on the universal cover
of $E_i$. Taking a logarithm of its fibre coordinate gives
the exact presentation
\[
 L_i^\times=\C^2/\Lambda_i,\qquad
 \Lambda_i=\left\langle
       (\gamma_1,\log\chi_i(\gamma_1)),
       (\gamma_2,\log\chi_i(\gamma_2)),(0,2\pi i)
                         \right\rangle_{\Z} .             \tag{C.MX.3}
\]
Here $\gamma_1,\gamma_2$ are an oriented integral period
basis of $E_i$, and $\chi_i$ is the chosen constant
factor character. The lattice is discrete of real rank
three and spans $\C^2$ over $\C$. Each marked-section
neighbourhood is the entire normal line bundle obtained
by removing the other section. Inverting the line
coordinate identifies the two descriptions on their
common complement. Their positive normal meridians are
opposite primitive elements of $\Lambda_i$.

For Hopf components use the actual entire normal
line-bundle neighbourhoods and logarithmic lattices
of (C.HC.4)--(C.HC.5), including their finite diagonal
deck groups. Every edge is the same explicit operation:
apply its elliptic base isomorphism, and send a nonzero
normal vector to its dual vector whose product is $t$.
It is a global biholomorphism of the two complements.
On universal covers it is affine. Indeed its base lift
is affine, and the derivative of its logarithmic normal
multiplier is doubly periodic and holomorphic, hence
constant. Thus it has a constant complex linear part
which maps $\Lambda_i$ isomorphically to $\Lambda_{i+1}$.
The product defines the actual integral automorphism
$M$ up to conjugacy. Changes of logarithms change only
translations by lattice vectors.

Every complement is smoothly $T^3\times\R$.
For (C.MX.3), a unitary flat metric on $L_i$ makes
the logarithm of its norm a linear height, with compact
three-torus level sets. The same description for a
Hopf complement was proved in UN46. Hence all component
cores and annular necks have the same rank-three
fundamental group, with collar maps exactly the
indicated integral maps. The finite normal-product
plumbing of UN46 applies without a changed local map:
choose disjoint section neighbourhoods, use $xy=t$
in each dual-line neck, and glue where one normal
norm exceeds a fixed inner radius. The two ends
are disjoint for sufficiently small $|t|$.
Closed intermediate cores and necks give properness;
the opposite open ends give the same Hausdorff
comparison as in that proof.

Van Kampen consequently gives
\[
                        \pi_1(Y_t)=\Lambda_1\rtimes_M\Z .
                                                        \tag{C.MX.4}
\]
The rank-three lattice injects in this presentation.
For an actual given smoothing, the divisor sections
and the product-preserving smooth tubular map used
in UN46 identify its collars with these very normal
products. Radial contraction of the tubular maps
gives the same maps on $\pi_1$. Thus (C.MX.4) holds
also for that smoothing. A two-torus has abelian
fundamental group, forcing $M=I$.

Conversely, when $M=I$, the complex linear cycle
product is the identity because $\Lambda_1$
spans $\C^2$. Constant changes of the two logarithmic
one-forms on successive pieces therefore glue them
to a closed holomorphic coframe on each compact
nearby fibre. The commuting dual fields have
complete complex flows and give the actual
biholomorphism with $\C^2$ divided by its period
lattice, exactly as in (C.HC.8). The core periods
and the one cycle period are all the fundamental
group generators, giving (C.MX.2). This proves
the complete local criterion for every stated
combination of components.

\par\medskip\noindent\textbf{Common normal periods and the Hopf restriction.}
Assume $M=I$ and identify all complements with
$G=\C^2/\Lambda$ using their actual linear maps.
Let $n_i\in\Lambda$ be the positive outgoing normal
meridian. The incoming normal meridian of component
$i$ is $-n_{i-1}$. On a ruled component the two
meridians are opposite, so
\[
                      n_i=n_{i-1}
                  \quad\hbox{at every ruled vertex}.   \tag{C.MX.5}
\]
At a Hopf vertex $n_i,n_{i-1}$ are complex-linearly
independent: its two logarithmic coordinate
meridians are a complex basis. A cycle with
exactly one Hopf vertex would have equality of
these two meridians by following (C.MX.5) through
all the other vertices. This proves its impossibility.

Put $W=\operatorname{span}_{\R}\Lambda$ and choose
a nonzero real linear functional $\ell$ with
kernel $W$. At a Hopf vertex, its contracting
logarithm gives
\[
 \log|\alpha_i|\,\ell(i n_i)
        =\log|\beta_i|\,\ell(i n_{i-1}).
                                                        \tag{C.MX.6}
\]
Both coefficients are negative, and neither value
can be zero because the two meridians span the
complex plane. At a ruled vertex their values
are equal by (C.MX.5). Thus, as soon as there is
a Hopf vertex, all $\ell(i n_i)$ have one nonzero
sign. The actual logarithmic cycle translation has
\[
                  2\pi i v_0=u=-\sum_i n_i,\qquad
                              \ell(iu)\ne0 .            \tag{C.MX.7}
\]
Write $u=d p$ with $p$ primitive in $\Lambda$ and
$d>0$. The projection along $\C p$ takes $\Lambda$
to a genuine elliptic lattice, since
$\C p\cap W=\R p$. Exactly as in (C.HG.6), actual
linear coordinates give
\[
 \Lambda=\left\langle
 (0,2\pi i),(2\pi i,0),(\log q_0,\log a)
                     \right\rangle_{\Z},\qquad
              |q_0|<1,\quad u=d(0,2\pi i).             \tag{C.MX.8}
\]
Here $a\ne1$: otherwise the explicit embedded
elliptic curve (C.HG.7) would lie in the common
complement of an algebraic-dimension-zero Hopf
surface, contradicting its two-curve calculation.

\par\medskip\noindent\textbf{Transporting both sorts of component.}
Use the compact elliptic surface and suspension
constructed in UN48 with the exact periods
\[
 \Lambda(t)+\Z(0,d\log t),\qquad
 \Lambda(t)=\left\langle
 (0,2\pi i),(2\pi i,0),(\log q_0,\log a(t))
                         \right\rangle_{\Z},             \tag{C.MX.9}
\]
where $a(0)=a$ and the actual Tate parameter is
$Q=t^d$. All the prescribed normal meridians and
contracting logarithms have fixed integral
coordinates in (C.MX.8). Retain those coordinates
in (C.MX.9).

At a Hopf vertex use the moving coordinate basis
$n_i(t)/(2\pi i),-n_{i-1}(t)/(2\pi i)$.
It is invertible near zero, retains both
contracting multipliers and all finite diagonal
characters, and gives the actual Hopf family
of UN48.

At a ruled vertex those two columns are
dependent, so a different actual frame is
required. Its primitive normal period $n_i$
extends to a lattice basis $n_i,\lambda_i,c_i$.
Projection along its complex line gives the
elliptic base lattice of this particular
ruled component. Choose a base coordinate
with the projected $\lambda_i$-period $2\pi i$,
and a logarithmic fibre coordinate having
normal period $2\pi i$ and zero $\lambda_i$-period.
Transport $n_i,\lambda_i,c_i$ by their fixed
integer coordinates. The first two remain
complex-linearly independent. They give
holomorphic coordinates in which the complement
lattice is
\[
        \left\langle
          (0,2\pi i),(2\pi i,0),
                      (\gamma_i(t),b_i(t))
                                  \right\rangle_{\Z} .   \tag{C.MX.10}
\]
Its projected periods remain real-linearly
independent on a smaller disc. They define
a proper elliptic family $E_i(t)$, and the
factor $\exp b_i(t)$ defines an actual
degree-zero line bundle $L_i(t)$ on it.
The resulting
$\PP(\OO_{E_i(t)}\oplus L_i(t))$ is a smooth
proper family of the prescribed ruled
components, with their marked sections.
At zero its line-bundle presentation is
the original one in (C.MX.3), so it is
an isomorphism of marked surfaces.

Every moving edge linear map sends the
outgoing normal vector to minus the next
incoming normal vector. Its induced base
map is therefore a holomorphic elliptic
isomorphism, and its nonzero normal map
is the reciprocal linear-fibre map with
the prescribed constant factors. Thus it
extends to exactly the required dual
line-bundle data. Keep the constants of
all original elliptic gluing maps.
The coefficients of the adjustable
normal-product logarithms in the total
common translation are
$-n_i(t)/(2\pi i)$. At a Hopf vertex
the two adjacent such vectors are
complex-linearly independent. Fix the
other constants and solve the resulting
two-by-two holomorphic system to make
the constant cycle translation zero.
The total period is then exactly
$(0,d\log t)$ in (C.MX.9), while every
central component and elliptic gluing
map is unchanged.

The uniformly chosen flat metrics, cores
and normal-product necks now give a
smooth proper local filling. On its
punctured disc the piecewise affine
maps into the common period quotient
agree modulo precisely (C.MX.9);
their changes of logarithm are integral
normal periods. Their fibrewise inverses
are the closed-coframe inverses already
proved above, so they give an actual
biholomorphism of the punctured
families with the compact suspension.
The nested-annulus replacement of
UN48 produces a smooth compact total
space over $\PP^1$ with exactly the
specified mixed central support.
Each of its components and double
curves has Euler number zero. All
other reduced supports are the
elliptic-base bundles in that
suspension and also have Euler
number zero. This proves the
compact assertion.

\par\medskip\noindent\textbf{The retained all-ruled translation.}
When every vertex is ruled, (C.MX.5)
makes all normal periods equal to
one primitive $n$. The common
elliptic quotient is
$\C^2/(\C n+\Lambda)$. The coefficient
vectors of all normal-product constants
lie in the single line $\C n$,
whereas the projection of the full
affine cycle constant $b_0$ is the
composition of the actual base
translations. Consequently varying
normal-product constants changes
only its vertical part; the period
formula (C.MX.2) retains the projected
translation. The local torus proof
above made no assumption that it
vanishes. In particular the
single ruled self-identification
is included by its actual
normal-line inversion, rather than
the Hopf coordinate interchange.
The exact examples and moving-frame comparisons in
\path{code/A_mixed_cycles.py}
check 169 mixed cycles with two Hopf vertices and
arbitrary numbers of ruled insertions in the
displayed range. Their primitive cycle divisor is
$\gcd(a,b)$, and their real period determinant is
$\log|t|\bigl(a\log|\beta|+b\log|\alpha|\bigr)>0$.
The same calculation retains the all-ruled
horizontal translation exactly.
\end{proof}


% C.2.UN.72 from part_C2_no_triple_classification.tex

